\section{Пересечение и сумма подпространств. Построение базисов в сумме и пересечении подпространств}

\begin{definition}
Пусть $U_1$ и $U_2$ --- подпространства линейного пространства $V$.
\begin{enumerate}
    \item \textbf{Суммой} подпространств $U_1$ и $U_2$ называется множество всех возможных сумм векторов из этих подпространств:
    \[
    U_1 + U_2 = \{ x \in V \mid x = x_1 + x_2, \; x_1 \in U_1, \; x_2 \in U_2 \}.
    \]
    \item \textbf{Пересечением} подпространств $U_1$ и $U_2$ называется множество векторов, принадлежащих обоим подпространствам одновременно:
    \[
    U_1 \cap U_2 = \{ x \in V \mid x \in U_1 \text{ и } x \in U_2 \}.
    \]
\end{enumerate}
Оба множества являются подпространствами в $V$.
\end{definition}

\begin{theorem}[О размерности суммы и пересечения]
Для любых конечномерных подпространств $U_1$ и $U_2$ пространства $V$ справедлива формула Грассмана:
\[
\dim(U_1 + U_2) = \dim U_1 + \dim U_2 - \dim(U_1 \cap U_2).
\]
\end{theorem}

\begin{proof}
Пусть $\dim(U_1 \cap U_2) = m$. Выберем базис пересечения: $f_1, \dots, f_m$. 
По теореме о неполном базисе, дополним эту систему до базиса в $U_1$: $f_1, \dots, f_m, e_1, \dots, e_k$ (где $\dim U_1 = m + k$), и до базиса в $U_2$: $f_1, \dots, f_m, g_1, \dots, g_l$ (где $\dim U_2 = m + l$).

Рассмотрим объединенную систему векторов:
\[
S = \{ f_1, \dots, f_m, e_1, \dots, e_k, g_1, \dots, g_l \}.
\]
\begin{enumerate}
    \item \textit{Система $S$ порождает $U_1 + U_2$:} Любой вектор $x \in U_1 + U_2$ представим как $x = x_1 + x_2$, где $x_1 \in U_1$, $x_2 \in U_2$. Так как $x_1$ линейно выражается через $f_i, e_j$, а $x_2$ --- через $f_i, g_p$, их сумма линейно выражается через всю систему $S$.
    \item \textit{Система $S$ линейно независима:} Составим линейную комбинацию, равную нулю:
    \[
    \sum_{i=1}^m \alpha_i f_i + \sum_{j=1}^k \beta_j e_j + \sum_{p=1}^l \gamma_p g_p = 0. \quad (*)
    \]
    Перенесем векторы $g_p$ в правую часть:
    \[
    \sum_{i=1}^m \alpha_i f_i + \sum_{j=1}^k \beta_j e_j = - \sum_{p=1}^l \gamma_p g_p.
    \]
    Левая часть равенства принадлежит $U_1$, а правая --- $U_2$. Следовательно, этот вектор принадлежит $U_1 \cap U_2$. Значит, он линейно выражается только через базис пересечения $f_1, \dots, f_m$:
    \[
    - \sum_{p=1}^l \gamma_p g_p = \sum_{i=1}^m \delta_i f_i \quad \Rightarrow \quad \sum_{i=1}^m \delta_i f_i + \sum_{p=1}^l \gamma_p g_p = 0.
    \]
    Так как система $\{f_1, \dots, f_m, g_1, \dots, g_l\}$ является базисом $U_2$, она линейно независима. Следовательно, все $\gamma_p = 0$ и все $\delta_i = 0$. 
    Подставляя $\gamma_p = 0$ в $(*)$, получаем $\sum \alpha_i f_i + \sum \beta_j e_j = 0$. В силу линейной независимости базиса $U_1$, все $\alpha_i = 0$ и $\beta_j = 0$.
\end{enumerate}
Таким образом, $S$ --- базис $U_1 + U_2$. Количество векторов в $S$ равно $m + k + l$. 
Заметим, что $\dim U_1 + \dim U_2 - \dim(U_1 \cap U_2) = (m+k) + (m+l) - m = m+k+l$. Теорема доказана.
\hfill$\blacksquare$
\end{proof}

\begin{remark}[Построение базисов на практике]
\begin{itemize}
    \item \textbf{Базис суммы $U_1 + U_2$:} объединяем базисы $U_1$ и $U_2$ в одну систему векторов, записываем их координаты в строки (или столбцы) матрицы и приводим её к ступенчатому виду. Ненулевые строки (или исходные векторы, соответствующие базисным столбцам) образуют базис суммы.
    \item \textbf{Базис пересечения $U_1 \cap U_2$:} если $U_1$ и $U_2$ заданы своими базисами, вектор $x \in U_1 \cap U_2$ должен линейно выражаться и через базис $U_1$, и через базис $U_2$. Приравнивая эти разложения, получаем однородную СЛАУ. Фундаментальная система решений (ФСР) этой СЛАУ дает коэффициенты для построения базиса пересечения.
\end{itemize}
\end{remark}
